\documentclass[11pt,a4paper]{article}

% --- PACKAGES BASE (COMPATIBLES TOUTES DISTRIBUTIONS LATEX) ---
\usepackage[utf8]{inputenc}
\usepackage[T1]{fontenc}
\usepackage{lmodern}
\usepackage{geometry}
\geometry{margin=2.5cm}
\usepackage{amsmath,amssymb,amsfonts}
\usepackage{graphicx}
\usepackage{xcolor}

% --- COULEURS PERSONNALISÉES ---
\definecolor{primaryblue}{RGB}{0, 51, 102}
\definecolor{secondarygreen}{RGB}{34, 139, 34}
\definecolor{darkgray}{RGB}{60, 60, 60}

\begin{document}

% ==============================================================================
% EN-TÊTE DU DOCUMENT / PAGE DE TITRE
% ==============================================================================
\begin{center}
    {\Large \textbf{ÉTABLISSEMENT D'ENSEIGNEMENT SUPÉRIEUR}}\\[0.2cm]
    {\large \textbf{Département de Génie Électrique}}\\[0.1cm]
    {\small \textbf{Projet de Fin d'Études / Projet d'Année -- Classe Terminale (Génie Électrique)}}\\[0.4cm]
    
    \rule{\linewidth}{0.8pt}\\[0.3cm]
    {\LARGE \bfseries \color{primaryblue} Déploiement Embarqué du Système Nerveux Bio-inspiré de la Drosophile sur Drone pour la Détection Prédictive d'Anomalies des Plantes} \\[0.3cm]
    \rule{\linewidth}{0.8pt}\\[0.4cm]
    
    \textbf{Domaines d'application :} Agriculture Intelligente et de Précision $\bullet$ IA Multimodale \& Traitement du Signal $\bullet$ Systèmes Embarqués \& Robotique \\
    \textbf{Présenté par :} Équipe d'Étudiants en Génie Électrique \\
    \textbf{Date de soumission du dossier :} Mercredi 7 Octobre 2026
\end{center}

\vspace{0.3cm}

\begin{abstract}
\noindent Ce document constitue le dossier de cadrage scientifique et technique pour notre projet d'année en Génie Électrique. Le projet vise à modéliser, simplifier et embarquer la dynamique du système nerveux (cerveau central CNS et chaîne nerveuse ventrale VNC) de la drosophile (\textit{Drosophila melanogaster}) sur un micro-drone d'inspection agricole. En combinant des capteurs d'olfaction chimique (détection précoce des composés organiques volatils COV) et une vision rétinienne bio-inspirée, le drone réalise un diagnostic prédictif des pathologies végétales avant l'apparition de symptômes visibles, tout en réduisant drastiquement la consommation énergétique embarquée grâce au calcul neuromorphique.
\end{abstract}

\vspace{0.2cm}

% ==============================================================================
% SECTION 1 : ÉTAT DE L'ART DU SUJET
% ==============================================================================
\section*{\color{primaryblue}1. État de l'Art du Sujet}
\hrule
\vspace{0.3cm}

\subsection*{\color{secondarygreen}1.1 Problématique en Agriculture Intelligente et de Précision}
L'agriculture de précision nécessite une surveillance continuous et automatisée de la santé des cultures afin de limiter les pertes d'exploitation et de cibler précisément l'application des traitements phytosanitaires.

Les approches de télédétection actuellement déployées par drones (UAV) présentent deux limites majeures :
\begin{itemize}
    \item \textbf{Diagnostic visuel tardif :} Les indices de végétation (ex. NDVI) et la classification par vision classique (RGB / multispectrale) n'identifient les foyers d'infection que lorsque le stress foliaire provoque des nécroses ou décolorations visibles à l'œil nu (phase avancée de la maladie).
    \item \textbf{Verrou énergétique embarqué :} Les algorithmes de vision par ordinateur reposant sur des réseaux de neurones profonds classiques (CNN) exigent une puissance de calcul élevée ($> 15\text{ à }30\text{ W}$), nécessitant des cartes d'accélération GPU lourdes et énergivores qui restreignent fortement l'autonomie de vol des micro-drones.
\end{itemize}

\subsection*{\color{secondarygreen}1.2 Le Système Nerveux de la Drosophile (\textit{Drosophila melanogaster}) comme Modèle Bio-inspiré}
La mouche du vinaigre (\textit{Drosophila melanogaster}) possède un système nerveux remarquablement compact ($\sim 10^5$ neurones) lui permettant de naviguer avec précision, d'éviter les obstacles et de localiser des sources d'odeurs à l'échelle du PPM (partie par million), le tout avec une consommation de quelques microwatts.

Les récents développements en neuro-ingénierie et en simulation bio-mécanique (notamment le framework \texttt{FlyGym} et le connectome \texttt{FlyLab}) ont mis en évidence la structure décisionnelle et sensori-motrice de l'insecte :
\begin{enumerate}
    \item \textbf{Sensibilité Olfactive Précoce (Antennes et Maxillaires) :} Lors d'une attaque parasitaire ou d'un stress abiotique, les plantes émettent immédiatement des composés organiques volatils (COV / \textit{VOCs}). Les sensilles olfactives de la drosophile capturent ces gradients chimiques très tôt.
    \item \textbf{Vision Composée Bio-inspirée :} Les ommatidies de l'œil composé filtrent les informations de contraste et de flux optique avec une latence inférieure à $10\text{ ms}$.
    \item \textbf{Contrôle Hiérarchique CNS--VNC :} Le cerveau (\textit{Central Nervous System}) effectue la fusion des données sensorielles et la prise de décision spatiale (anémotaxie), tandis que la chaîne nerveuse ventrale (VNC -- \textit{Ventral Nerve Cord}) génère les commandes réflexes et les motifs d'action sensori-moteurs.
\end{enumerate}

\subsection*{\color{secondarygreen}1.3 Neuromorphic Computing et Systèmes Embarqués pour Drones}
Le transfert de ces modèles biologiques vers des cibles microélectroniques s'appuie sur le calcul neuromorphique et les réseaux de neurones à impulsions (\textit{Spiking Neural Networks} -- SNN). Dans ce paradigme événementiel (\textit{event-driven}), les informations ne sont traitées que lors de l'émission d'impulsions (\textit{spikes}), réduisant la consommation d'énergie statique et dynamique d'un facteur 100 par rapport aux microprocesseurs classiques.

\subsection*{\color{secondarygreen}1.4 Intérêt et Importance Scientifique et Technique du Projet}
Pour un futur ingénieur en Génie Électrique, ce projet revêt une importance majeure à plusieurs niveaux :
\begin{itemize}
    \item \textbf{Innovation technologique :} Fusion de la robotique aérienne autonome, du calcul neuromorphique basse consommation et des capteurs physico-chimiques.
    \item \textbf{Impact socio-économique :} Permettre une détection prédictive (au stade pré-symptomatique) des maladies des plantes pour préserver la sécurité alimentaire et limiter les intrants chimiques.
    \item \textbf{Synergie disciplinaire :} Application concrète de l'électronique embarquée, du traitement du signal/image, de l'automatique et de l'intelligence artificielle multimodale.
\end{itemize}

% ==============================================================================
% SECTION 2 : CAHIER DES CHARGES PRATIQUE
% ==============================================================================
\vspace{0.4cm}
\section*{\color{primaryblue}2. Cahier des Charges Pratique}
\hrule
\vspace{0.3cm}

\subsection*{\color{secondarygreen}2.1 Description des Objectifs du Projet}
L'objectif principal du projet est la conception, la modélisation, le prototypage électronique et la validation d'une architecture de contrôle et de diagnostic bio-inspirée embarquée sur micro-drone.

\subsubsection*{\color{darkgray}2.1.1 Objectifs Spécifiques :}
\begin{enumerate}
    \item \textbf{Modélisation \& Conversion Neuromorphique :} Extraire la topologie sensori-motrice du système nerveux de la drosophile depuis l'environnement \texttt{FlyGym}/\texttt{FlyLab} et la convertir en un modèle SNN compact.
    \item \textbf{Conception Électronique Embarquée :} Concevoir une carte de traitement et d'acquisition électronique optimisée autour d'un système embarqué rapide (ex. STM32H7 / ESP32-S3 / NVIDIA Jetson Nano / NPU Coral).
    \item \textbf{Banc d'Acquisition Sensoriel Multimodal :} Développer et conditionner les signaux de capteurs d'olfaction chimique (nez électronique MOX/COV) et d'un capteur optique/caméra miniature.
    \item \textbf{Implémentation des Algorithmes de Vol \& Diagnostic :} Programmer en C/C++ l'algorithme d'orientation vers la source d'odeur (anémotaxie bio-inspirée) et le module de diagnostic prédictif des anomalies.
    \item \textbf{Validation et Benchmarking :} Évaluer expérimentalement le temps de réponse, la précision du diagnostic et l'économie d'énergie réalisée ($< 3\text{ W}$).
\end{enumerate}

\subsection*{\color{secondarygreen}2.2 Description des Tâches à Réaliser (Découpage du Projet / WBS)}

\begin{description}
    \item[Tâche 1 : Étude Bibliographique et Prise en Main (Semaines 1--2)] \hfill \\
    Analyse théorique du modèle \texttt{FlyGym}/MuJoCo, sélection des capteurs de gaz COV et prise en main des outils d'apprentissage neuromorphique.
    
    \item[Tâche 2 : Modélisation et Compression du Réseau Nerveux (Semaines 3--4)] \hfill \\
    Simplification de la dynamique neuro-fonctionnelle de la drosophile pour créer un réseau SNN à faible latence optimisé pour microcontrôleur.
    
    \item[Tâche 3 : Conception Électronique et Hardware (Semaines 5--6)] \hfill \\
    Dimensionnement du circuit d'alimentation, routage PCB/maquette d'interface capteurs, intégration matérielle sur le châssis du drone.
    
    \item[Tâche 4 : Programmation et Intégration du Firmware (Semaines 7--8)] \hfill \\
    Développement du code embarqué C/C++ pour l'acquisition temps réel des capteurs, l'exécution du réseau SNN et la génération des consignes de vol.
    
    \item[Tâche 5 : Tests, Validation Expérimentale et Rapport (Semaines 9--10)] \hfill \\
    Essais en environnement contrôlé (émission de gaz COV sur plantes), mesure de la consommation électrique, optimisation finale et rédaction du rapport.
\end{description}

\subsection*{\color{secondarygreen}2.3 Moyens Nécessaires à la Mise en Œuvre du Projet}

\subsubsection*{\color{darkgray}2.3.1 Moyens Logiciels (\textit{Software}) :}
\begin{itemize}
    \item \textbf{Simulation Bio-mécanique \& Neuro-ingénierie :} Python 3.10+, \texttt{FlyGym}, \texttt{MuJoCo}.
    \item \textbf{Apprentissage \& Neuromorphisme :} PyTorch, \texttt{Brian2} / \texttt{SNNorch}, TensorFlow Lite.
    \item \textbf{Développement Embarqué \& CAO Électronique :} STM32CubeIDE / VS Code PlatformIO, Altium Designer / KiCad.
\end{itemize}

\subsubsection*{\color{darkgray}2.3.2 Moyens Matériels (\textit{Hardware}) :}
\begin{itemize}
    \item \textbf{Unité de Traitement Embarquée :} Carte de calcul STM32H753 / ESP32-S3 ou Raspberry Pi 4 avec accélérateur Google Coral NPU.
    \item \textbf{Capteurs d'Olfaction (Nez Électronique) :} Matrice de capteurs de gaz MOX (BME688 / SGP40 / MQ-135) sensibles aux COV légers (éthylène, jasmonate, éthanol).
    \item \textbf{Capteurs de Navigation \& Vision :} Caméra OV5640 / ArduCam miniature, capteur de flux optique (PMW3901) et IMU (MPU6050 / ICM-20948).
    \item \textbf{Vecteur Aérien :} Châssis de drone quadricoptère compact 2--3 pouces avec contrôleur de vol commercial (Betaflight / PX4).
    \item \textbf{Banc d'Essai d'Électronique :} Alimentation stabilisée de laboratoire, oscilloscope numérique, multimètres et station de soudage.
\end{itemize}

% ==============================================================================
% PLANNING PRÉVISIONNEL
% ==============================================================================
\subsection*{\color{secondarygreen}2.4 Planning Prévisionnel et Jalons du Projet}

Le tableau 1 synthétise le calendrier prévisionnel des travaux jusqu'à la livraison finale le 30 novembre 2026.

\vspace{0.2cm}
\begin{center}
\small
\begin{tabular}{|l|p{6.5cm}|c|p{4cm}|}
\hline
\textbf{Phase} & \textbf{Intitulé de la Tâche} & \textbf{Période} & \textbf{Livrable associé} \\ \hline
Phase 1 & Cadrage, État de l'art \& Cahier des charges & 01/10 -- 07/10/2026 & Dossier de cadre (5 pages) \\ \hline
Phase 2 & Modélisation du Système Nerveux \& SNN & 08/10 -- 20/10/2026 & Code SNN \& Validation MuJoCo \\ \hline
Phase 3 & Conception Matérielle \& Cartes Interface & 21/10 -- 05/11/2026 & Schéma PCB \& Proto matériel \\ \hline
Phase 4 & Dével. Firmware Embarqué C/C++ & 06/11 -- 20/11/2026 & Firmware sur carte cible \\ \hline
Phase 5 & Tests, Validation \& Soutenance & 21/11 -- 30/11/2026 & Drone fonctionnel \& Rapport final \\ \hline
\end{tabular}
\end{center}
\vspace{0.2cm}

% ==============================================================================
% CONCLUSION ET PERSPECTIVES
% ==============================================================================
\section*{\color{primaryblue}3. Conclusion et Perspectives}
\hrule
\vspace{0.3cm}
Ce projet d'année apporte une contribution originale et pertinente au domaine de l'agriculture intelligente. En s'inspirant des facultés exceptionnelles du système nerveux de la drosophile, la solution proposée permet d'associer la détection prédictive chimique (COV) et la frugalité énergétique du calcul neuromorphique. Les compétences développées couvrent l'ensemble du programme de Génie Électrique : électronique embarquée, automatique, traitement du signal bio-inspiré et robotique aérienne.

\end{document}
